%%%%%%%%%%%%%%%%%%%%%%%%%%%%%%%%%%%%%%%%%%%%%%%%%%%%%%%%%%%%%%%%%%%%%%%%%%%%%%%%%%%%
% F.CS308 Машин сургалт — Лекц 5: Neural Networks
% Лекц дээр тавигдсан бүх бодлогын бодолт
%%%%%%%%%%%%%%%%%%%%%%%%%%%%%%%%%%%%%%%%%%%%%%%%%%%%%%%%%%%%%%%%%%%%%%%%%%%%%%%%%%%%
\documentclass[11pt]{article}
\usepackage[utf8]{inputenc}            % UTF8 юникод текст оруулах
\usepackage[T2A]{fontenc}              % кирил үсгийн кодчилол
\usepackage[mongolian]{babel}          % олон хэлний текст
\usepackage[margin=1in]{geometry}
\usepackage{amsmath,amsthm,amssymb,amsfonts,fancyhdr,color,comment,graphicx,environ}
\usepackage{xcolor}
\usepackage{mdframed}
\usepackage[shortlabels]{enumitem}
\usepackage{indentfirst}
\usepackage{longtable}
\usepackage[table,xcdraw]{xcolor}
\usepackage{listings}
\usepackage{booktabs}
\usepackage{mathtools}                 % psmallmatrix
\usepackage{pifont}                    % \ding{55}
\usepackage{hyperref}
\hypersetup{colorlinks=true, linkcolor=blue, filecolor=magenta, urlcolor=blue}

\pagestyle{fancy}
\lhead{B231910012, Б. Өсөхбаяр \\ F.CS308 Машин сургалт, Лекц 5 — Даалгавар}
\rhead{\thepage}
\setlength{\headheight}{26pt}

\newenvironment{problem}[2][Бодлого]
    { \begin{mdframed}[backgroundcolor=gray!20] \textbf{#1 #2} \\}
    { \end{mdframed}}

\newenvironment{solution}[1][\textbf{Бодолт.}]
    { \noindent #1 \smallskip \\ }
    { \medskip }

\newcommand{\xin}{\mathbf{x}_{\mathrm{in}}}
\newcommand{\xout}{\mathbf{x}_{\mathrm{out}}}
\newcommand{\dL}[1]{\frac{\partial \mathcal{L}}{\partial #1}}

\newenvironment{codebox}
  {\begin{mdframed}[backgroundcolor=gray!8,linewidth=0.4pt]\footnotesize}
  {\end{mdframed}}

\title{\textbf{Лекц 5: Neural Networks} \\[4pt]
       \large Лекц дээр тавигдсан бүх бодлогын бодолт}
\author{Б. Өсөхбаяр \quad (B231910012)}
\date{2026 оны 10 сарын 15}

\begin{document}
\maketitle

\noindent
Энэхүү тайланд \textbf{F.CS308 Лекц 5 (Neural Networks, 147 слайд)}-д тавигдсан бүх
бодлогыг бодов. Бодлого бүрийн ард лекцийн тухайн слайдын дугаарыг бичсэн болно.
Бүх тоон үр дүнг NumPy дээр бие даан шалгасан (Хавсралт А).

\tableofcontents
\bigskip
\hrule

%==================================================================================
\section{Бодлого 1 — Гинжний дүрмийн матрицын хэмжээсийг шалгах}
%==================================================================================

\begin{problem}{1 (слайд 115)}
Дараах гинжин уламжлалын \textbf{бүх матрицын хэмжээс} таарч байгаа эсэхийг шалга:
\[
\left.\frac{\partial \mathbf{z}}{\partial \mathbf{x}}\right|_{\mathbf{x}=\mathbf{a}}
= \left.\frac{\partial \mathbf{z}}{\partial \mathbf{u}_{L-1}}\right|
\cdot \left.\frac{\partial \mathbf{u}^{L-1}}{\partial \mathbf{u}_{L-2}}\right|
\cdots
\left.\frac{\partial \mathbf{u}_{2}}{\partial \mathbf{u}_{1}}\right|
\cdot \left.\frac{\partial \mathbf{u}_{1}}{\partial \mathbf{x}}\right|_{\mathbf{x}=\mathbf{a}}
\]
энд $h(\mathbf{x}) = f_L(f_{L-1}(\dots f_1(\mathbf{x})))$,
$\mathbf{u}_1 = f_1(\mathbf{x})$, $\mathbf{u}_l = f_l(\mathbf{u}_{l-1})$,
$\mathbf{z} = \mathbf{u}_L = f_L(\mathbf{u}_{L-1})$.
\end{problem}

\begin{solution}
\textbf{Алхам 1. Jacobian-ы хэмжээсийн дүрмийг тогтооё.}
Лекцийн ``Matrix тооцоо'' слайдаар:
\begin{center}
\begin{tabular}{lll}
\toprule
\textbf{Юуны уламжлал} & \textbf{Тэмдэглэгээ} & \textbf{Үр дүнгийн хэмжээс} \\
\midrule
скаляр / вектор $[n\times 1]$       & $\partial y / \partial \mathbf{x}$        & мөр вектор $[1\times n]$ \\
вектор $[m\times 1]$ / вектор $[n\times 1]$ & $\partial \mathbf{y} / \partial \mathbf{x}$ & матриц $[m\times n]$ \\
скаляр / матриц $[n\times m]$       & $\partial y / \partial \mathbf{X}$        & матриц $[m\times n]$ \\
\bottomrule
\end{tabular}
\end{center}

\noindent
Тэмдэглэе: $n = |\mathbf{x}|$, \quad $p_l = |\mathbf{u}_l|$, \quad
$m = |\mathbf{z}| = p_L$.

\medskip
\textbf{Алхам 2. Үржигдэхүүн бүрийн хэмжээсийг бичье.}
\begin{align*}
\frac{\partial \mathbf{z}}{\partial \mathbf{u}_{L-1}}         &: \quad [\,m \times p_{L-1}\,] \\
\frac{\partial \mathbf{u}_{L-1}}{\partial \mathbf{u}_{L-2}}   &: \quad [\,p_{L-1} \times p_{L-2}\,] \\
&\;\;\vdots \\
\frac{\partial \mathbf{u}_{2}}{\partial \mathbf{u}_{1}}       &: \quad [\,p_{2} \times p_{1}\,] \\
\frac{\partial \mathbf{u}_{1}}{\partial \mathbf{x}}           &: \quad [\,p_{1} \times n\,]
\end{align*}

\textbf{Алхам 3. Үржвэрийг гүйцэтгэе.}
Матриц үржихэд зэргэлдээ хэмжээсүүд таарах ёстой. Дээрх дараалалд:
\[
\underbrace{[m \times \mathbf{p_{L-1}}]}_{}\;
\underbrace{[\mathbf{p_{L-1}} \times \mathbf{p_{L-2}}]}_{}\;
\cdots\;
\underbrace{[\mathbf{p_{2}} \times \mathbf{p_{1}}]}_{}\;
\underbrace{[\mathbf{p_{1}} \times n]}_{}
\;=\; [\,m \times n\,]
\]
Дотоод хэмжээс бүр ($p_{L-1}, p_{L-2}, \dots, p_1$) \textbf{яг хоёр удаа} давхцаж
хоорондоо цуцлагдаж байна. Үлдэх нь хамгийн зүүн талын $m$ ба хамгийн баруун талын $n$.

\medskip
\textbf{Алхам 4. Хүлээгдэж буй хэмжээстэй харьцуулъя.}
$\mathbf{z} \in \mathbb{R}^m$, $\mathbf{x} \in \mathbb{R}^n$ тул дүрмийн 2-р мөрөөр
$\partial \mathbf{z} / \partial \mathbf{x}$ нь $[m \times n]$ байх ёстой.
Энэ нь Алхам 3-ын үр дүнтэй \textbf{яг таарч байна}. $\blacksquare$

\medskip
\textbf{Тоон жишээ (лекцийн өөрийнх нь жишээ).}
$|\mathbf{z}| = 1$, $|\mathbf{u}| = 2$, $|\mathbf{x}| = 4$ үед:
\[
h'(\mathbf{x}) \;=\; \underbrace{[1 \times 2]}_{\partial \mathbf{z}/\partial \mathbf{u}}
\;\cdot\; \underbrace{[2 \times 4]}_{\partial \mathbf{u}/\partial \mathbf{x}}
\;=\; [1 \times 4] \quad \checkmark
\]
Энэ нь лекцийн слайд дээрх цэнхэр ($1\times2$) ба улаан ($2\times4$) блокуудын зураглалтай
яг тохирч байна.

\medskip
\begin{mdframed}[backgroundcolor=yellow!12]
\textbf{Чухал дүгнэлт.} $\mathbf{z}$ нь \emph{скаляр} байх тохиолдолд ($m=1$) үр дүн нь
\textbf{мөр вектор} $[1 \times n]$ болно. Яг энэ шалтгаанаар алдагдлын давхаргын градиент
\[
\dL{\mathbf{x}_L} = (\mathbf{x}_L - \mathbf{y})^{T}
\]
гэж \textbf{транспозтой} бичигддэг. Хэрэв транспозыг орхивол дараагийн алхмын
$\dL{\mathbf{x}_{L-1}} = \dL{\mathbf{x}_L}\cdot \mathbf{W}_L$ үржвэрийн хэмжээс
таарахгүй болно.
\end{mdframed}
\end{solution}

%==================================================================================
\section{Бодлого 2 — XOR бодлого}
%==================================================================================

\begin{problem}{2 (слайд 9)}
XOR функц:
\begin{center}
\begin{tabular}{ccc}
\toprule
$x_1$ & $x_2$ & гаралт \\
\midrule
0 & 0 & 0 \\
1 & 0 & 1 \\
0 & 1 & 1 \\
1 & 1 & 0 \\
\bottomrule
\end{tabular}
\end{center}
(а) Нэг давхаргат перцептрон үүнийг \textbf{бодож чадахгүйг} батал. \\
(б) Хоёр давхаргат сүлжээгээр \textbf{тодорхой жин бичиж} шийд.
\end{problem}

\begin{solution}[\textbf{(а) Боломжгүйн баталгаа.}]
Эсрэгээс нь тооцъё: $\exists\, \mathbf{w}=(w_1,w_2)$, $b$ байж
$f(\mathbf{x}) = \mathbb{1}(w_1x_1 + w_2x_2 + b > 0)$ нь XOR-ыг гаргадаг гэе.
Дөрвөн мөр тус бүрээс:
\begin{align}
f(0,0) = 0 &\;\Longrightarrow\; b \le 0 \\
f(1,0) = 1 &\;\Longrightarrow\; w_1 + b > 0 \\
f(0,1) = 1 &\;\Longrightarrow\; w_2 + b > 0 \\
f(1,1) = 0 &\;\Longrightarrow\; w_1 + w_2 + b \le 0
\end{align}

(2) ба (3)-ыг нэмбэл:
\[
w_1 + w_2 + 2b > 0 \quad\Longrightarrow\quad w_1 + w_2 > -2b
\]
Харин (4)-өөс:
\[
w_1 + w_2 \le -b
\]
Хоёуланг нь нийлүүлбэл
\[
-2b \;<\; w_1 + w_2 \;\le\; -b
\]
Энэ завсар хоосон биш байхын тулд $-2b < -b$, өөрөөр хэлбэл $\boxed{b > 0}$ байх ёстой.
Гэтэл (1)-ээс $b \le 0$. \textbf{Зөрчил.} $\blacksquare$

\medskip
\textbf{Геометр тайлбар.} $(0,0)$ ба $(1,1)$ цэгүүдийн дундаж нь $(0.5,\,0.5)$;
$(1,0)$ ба $(0,1)$ цэгүүдийн дундаж нь мөн $(0.5,\,0.5)$. Шугаман функц нь дундажийг
хадгалдаг тул энэ нэг цэг дээр хоёр ангийн дундаж утга \emph{ижил} гарах ёстой —
иймээс тэднийг ялгах шулуун оршин байх боломжгүй.
\end{solution}

\begin{solution}[\textbf{(б) Хоёр давхаргат шийдэл.}]
\textbf{Хувилбар 1 — step идэвхжүүлэгчээр.}
XOR $=$ (OR) AND (NOT AND) гэдгийг ашиглая.
\[
\mathbf{W}_1 = \begin{pmatrix} 1 & 1 \\ 1 & 1 \end{pmatrix},\quad
\mathbf{b}_1 = \begin{pmatrix} -0.5 \\ -1.5 \end{pmatrix},\quad
\mathbf{W}_2 = \begin{pmatrix} 1 & -1 \end{pmatrix},\quad
b_2 = -0.5
\]
\[
\mathbf{h} = \mathrm{step}(\mathbf{W}_1\mathbf{x} + \mathbf{b}_1),
\qquad y = \mathrm{step}(\mathbf{W}_2\mathbf{h} + b_2)
\]
$h_1$ нь \textbf{OR}, $h_2$ нь \textbf{AND}-ыг тооцоолно. Шалгая:
\begin{center}
\begin{tabular}{cc|cc|c|c}
\toprule
$x_1$ & $x_2$ & $h_1$ (OR) & $h_2$ (AND) & $y$ & XOR \\
\midrule
0 & 0 & $\mathrm{step}(-0.5)=0$ & $\mathrm{step}(-1.5)=0$ & $\mathrm{step}(-0.5)=0$ & 0 \checkmark \\
1 & 0 & $\mathrm{step}(0.5)=1$  & $\mathrm{step}(-0.5)=0$ & $\mathrm{step}(0.5)=1$  & 1 \checkmark \\
0 & 1 & $\mathrm{step}(0.5)=1$  & $\mathrm{step}(-0.5)=0$ & $\mathrm{step}(0.5)=1$  & 1 \checkmark \\
1 & 1 & $\mathrm{step}(1.5)=1$  & $\mathrm{step}(0.5)=1$  & $\mathrm{step}(-0.5)=0$ & 0 \checkmark \\
\bottomrule
\end{tabular}
\end{center}

\medskip
\textbf{Хувилбар 2 — ReLU идэвхжүүлэгчээр} (практикт ашиглагддаг хэлбэр):
\[
\mathbf{W}_1 = \begin{pmatrix} 1 & 1 \\ 1 & 1 \end{pmatrix},\quad
\mathbf{b}_1 = \begin{pmatrix} 0 \\ -1 \end{pmatrix},\quad
\mathbf{W}_2 = \begin{pmatrix} 1 & -2 \end{pmatrix},\quad
b_2 = 0
\]
\begin{center}
\begin{tabular}{cc|c|c|c}
\toprule
$x_1$ & $x_2$ & $\mathbf{h} = \mathrm{ReLU}(\mathbf{W}_1\mathbf{x}+\mathbf{b}_1)$ & $y = \mathbf{W}_2\mathbf{h}$ & XOR \\
\midrule
0 & 0 & $(0,\,0)$ & $0 - 0 = 0$ & 0 \checkmark \\
1 & 0 & $(1,\,0)$ & $1 - 0 = 1$ & 1 \checkmark \\
0 & 1 & $(1,\,0)$ & $1 - 0 = 1$ & 1 \checkmark \\
1 & 1 & $(2,\,1)$ & $2 - 2 = 0$ & 0 \checkmark \\
\bottomrule
\end{tabular}
\end{center}

\medskip
\begin{mdframed}[backgroundcolor=green!10]
\textbf{Яагаад ажиллав?} Нуугдмал давхарга $(x_1,x_2)$ орон зайг $(h_1,h_2)$ орон зай
руу хувиргав. Шинэ орон зайд дөрвөн цэг $(0,0), (1,0), (1,0), (2,1)$ болж — $(1,0)$ цэг
дээр \textbf{хоёр жишээ давхцаж}, үлдсэн хоёр нь шулуунаар тусгаарлагдах болсон.
Энэ нь \textbf{Лаб 4-ийн kernel trick}-тэй ижил санаа; ялгаа нь kernel-ийг гараар
сонгодог бол энд сүлжээ хувиргалтаа \emph{өөрөө сурдаг}.
\end{mdframed}
\end{solution}

%==================================================================================
\section{Бодлого 3 — Convolution module}
%==================================================================================

\begin{problem}{3 (слайд 123, гэрийн даалгавар)}
Оролт $\xin$ болон гаралт $\xout$ нь ижил урт $N$-ийн 1D дохио. Конволюцийн кернел нь
$\mathbf{W}$ бөгөөд урт $K < N$.
\textbf{Хайрцаг бүрийн доторх тэгшитгэлүүдийг гаргаж ав:} урагшлах тархалт,
буцах тархалт (оролт руу), жингийн градиент. Хил хязгаарын тухай ярилц.
\end{problem}

\begin{solution}
\subsection*{3.1 Урагшлах тархалт (ногоон хайрцаг)}
Корреляцийн хэлбэрээр тодорхойлъё (машин сургалтад түгээмэл конвенц):
\begin{equation}
\boxed{\;\xout[i] \;=\; \sum_{k=0}^{K-1} \mathbf{W}[k]\;\xin[i+k], \qquad i = 0,\dots,N-1\;}
\end{equation}
Гаралт нь оролттой ижил урттай байхын тулд $\xin[j] = 0$ для $j \ge N$ гэж
\textbf{тэгээр дүүргэнэ} (zero padding).

\subsection*{3.2 Буцах тархалт — оролт руу (улаан хайрцаг)}
Эхлээд нэг элементийн хэсэгчилсэн уламжлалыг олъё:
\[
\frac{\partial\, \xout[i]}{\partial\, \xin[j]} =
\begin{cases}
\mathbf{W}[j-i], & 0 \le j-i \le K-1 \\[2pt]
0, & \text{бусад тохиолдолд}
\end{cases}
\]
Гинжний дүрмээр, $\xin[j]$ нь \textbf{олон} гаралтад нөлөөлдөг тул бүгдийг нь нэмнэ
(лекцийн \emph{branch} дүрэм):
\begin{equation}
\boxed{\;\dL{\xin[j]} \;=\; \sum_{i=\max(0,\,j-K+1)}^{j} \dL{\xout[i]}\;\mathbf{W}[j-i]\;}
\end{equation}
Энэ нь $\dL{\xout}$ ба $\mathbf{W}$-ийн \textbf{конволюци} (индекс $j-i$) —
өөрөөр хэлбэл \emph{эргүүлсэн} кернелтэй корреляци. Урагшлах нь корреляци, буцах нь
конволюци болж байгааг анзаар.

\subsection*{3.3 Жингийн градиент (хоёр дахь улаан хайрцаг)}
\[
\frac{\partial\, \xout[i]}{\partial\, \mathbf{W}[k]} = \xin[i+k]
\]
$\mathbf{W}[k]$ нь \textbf{бүх} байрлалд дахин ашиглагддаг тул бүх $i$-ээр нэмнэ:
\begin{equation}
\boxed{\;\dL{\mathbf{W}[k]} \;=\; \sum_{i=0}^{N-1} \dL{\xout[i]}\;\xin[i+k],
\qquad k = 0,\dots,K-1\;}
\end{equation}

\begin{mdframed}[backgroundcolor=yellow!12]
\textbf{Linear module-тай гол ялгаа.} Linear давхаргад жин бүр \emph{нэг} л удаа
оролцдог тул $\dL{\mathbf{W}} = \xin \cdot \dL{\xout}$ (нийлбэргүй). Конволюцид жин
\textbf{хуваалцагддаг} (weight sharing) тул $N$ байрлалаас ирсэн бүх градиентийг
\emph{нэмж} хуримтлуулна. Энэ нь лекцийн ``Branching'' слайдын дүрмийн шууд үр дагавар:
\emph{нэг утга олон салаа руу явбал, буцахдаа градиентүүд нэмэгдэнэ}.
\end{mdframed}

\subsection*{3.4 Матрицын хэлбэр — Linear module-той холбоо}
Конволюци бол \textbf{Toeplitz} (туузан) матрицаар үржүүлэх явдал:
\[
\mathbf{T}[i,j] =
\begin{cases} \mathbf{W}[j-i], & 0 \le j-i < K \\ 0, & \text{бусад}\end{cases}
\qquad\Longrightarrow\qquad \xout = \mathbf{T}\,\xin
\]
Жишээ нь $N=4$, $K=3$ үед:
\[
\mathbf{T} = \begin{pmatrix}
W_0 & W_1 & W_2 & 0   \\
0   & W_0 & W_1 & W_2 \\
0   & 0   & W_0 & W_1 \\
0   & 0   & 0   & W_0
\end{pmatrix}
\]
Ийнхүү Linear module-ийн томъёо $\dL{\xin} = \dL{\xout}\cdot\mathbf{T}$ шууд хэрэглэгдэнэ.
Конволюци бол \emph{тусгай бүтэцтэй шугаман давхарга} л гэсэн үг.

\subsection*{3.5 Жингийн шинэчлэлт}
\[
\mathbf{W}^{k+1} \;\leftarrow\; \mathbf{W}^{k} - \eta \left(\frac{\partial J}{\partial \mathbf{W}}\right)^{T},
\qquad J = \sum_{i=1}^{N_{\text{жишээ}}} \mathcal{L}^{(i)}
\]

\subsection*{3.6 Тоон жишээ}
$N=4$, $K=3$, $\xin = (1,2,3,4)$, $\mathbf{W} = (1,0,-1)$, тэг дүүргэлттэй.

\textbf{Forward:}
\begin{align*}
\xout[0] &= 1\cdot1 + 0\cdot2 + (-1)\cdot3 = -2 \\
\xout[1] &= 1\cdot2 + 0\cdot3 + (-1)\cdot4 = -2 \\
\xout[2] &= 1\cdot3 + 0\cdot4 + (-1)\cdot0 = \phantom{-}3 \\
\xout[3] &= 1\cdot4 + 0\cdot0 + (-1)\cdot0 = \phantom{-}4
\end{align*}
$\Longrightarrow \xout = (-2,\,-2,\,3,\,4)$

\medskip
\textbf{Backward}, $\dL{\xout} = (1,1,1,1)$ гэж авбал:
\begin{align*}
\dL{\mathbf{W}[0]} &= 1+2+3+4 = 10 \\
\dL{\mathbf{W}[1]} &= 2+3+4+0 = \phantom{1}9 \\
\dL{\mathbf{W}[2]} &= 3+4+0+0 = \phantom{1}7
\end{align*}
\begin{align*}
\dL{\xin[0]} &= 1\cdot W_0 = 1 \\
\dL{\xin[1]} &= 1\cdot W_1 + 1\cdot W_0 = 0 + 1 = 1 \\
\dL{\xin[2]} &= 1\cdot W_2 + 1\cdot W_1 + 1\cdot W_0 = -1+0+1 = 0 \\
\dL{\xin[3]} &= 1\cdot W_2 + 1\cdot W_1 + 1\cdot W_0 = -1+0+1 = 0
\end{align*}
$\Longrightarrow \dL{\xin} = (1,\,1,\,0,\,0)$

\subsection*{3.7 Хил хязгаарын тухай хэлэлцүүлэг}
Бодлогод ``$\xin$ ба $\xout$ нь ижил урттай'' гэж заасан. $K > 1$ үед энэ нь
\textbf{заавал дүүргэлт (padding) шаардана}:
\begin{itemize}[leftmargin=*]
  \item \textbf{Тэг дүүргэлт (zero padding).} $K-1$ ширхэг тэг нэмнэ. Хамгийн энгийн,
        гэхдээ захын гаралтууд \emph{зохиомол} тэгүүдийг уншдаг тул бага зэрэг хазайна.
        Тэг нь нийлбэрийн хувьд \emph{нейтраль элемент} ($a+0=a$) учраас аюулгүй.
  \item \textbf{Толин тусгал / давтах (reflect / replicate).} Зургийн боловсруулалтад
        захын артефактыг багасгахад илүү тохиромжтой.
  \item \textbf{``Valid'' (дүүргэлтгүй).} Гаралтын урт $N-K+1 < N$ болно —
        бодлогын нөхцөлийг \emph{зөрчинө}, тиймээс энд тохирохгүй.
\end{itemize}
Буцах тархалтын үед: дүүргэсэн байрлалаас уншсан гишүүд $\dL{\mathbf{W}}$-д \textbf{0}
хувь нэмэр оруулна (учир нь $\xin = 0$); дүүргэлтийн байрлалд хуримтлагдсан
$\dL{\xin}$-ийг \textbf{хаяна} (тэд жинхэнэ оролт биш).

\medskip
\textbf{Шийдвэрийн заагийн (decision boundary) тухай.} Конволюци бол хуваалцсан жинтэй
\emph{шугаман} оператор. Иймээс зөвхөн конволюци давхаргууд дараалвал (шугаман бус
функцгүйгээр) нийлбэр үр дүн нь мөн л \textbf{шугаман} хэвээр үлдэж, шийдвэрийн зааг нь
гипер хавтгай л байна — Бодлого 2-ын XOR-ыг бодож чадахгүй. Шийдвэрийн заагийг
шугаман бус болгодог нь \textbf{ReLU ба pooling}. Конволюцийн өгдөг давуу тал нь
шугаман бус байдал биш, харин \emph{шилжилтийн эквивариант} ($\xin$-ийг шилжүүлбэл
$\xout$ нь ижил хэмжээгээр шилжинэ) ба параметрийн тоо эрс багасах явдал юм:
бүрэн холбоот давхаргад $N^2$ параметр хэрэгтэй бол конволюцид ердөө $K$.
\end{solution}

%==================================================================================
\section{Бодлого 4 — Max pooling module}
%==================================================================================

\begin{problem}{4 (слайд 124, гэрийн даалгавар)}
Оролтын $\xin$ болон гаралт $\xout$ нь \textbf{өөр өөр} урттай 1D дохио.
Хайрцаг бүрийн доторх тэгшитгэлүүдийг гаргаж ав.
Хил хязгаарыг хэрхэн даван туулах тухайгаа ярилц.
\end{problem}

\begin{solution}
\subsection*{4.1 Тэмдэглэгээ}
Цонхны хэмжээ $P$, алхам (stride) $S$, оролтын урт $N$. Гаралтын урт:
\[
M = \left\lfloor \frac{N-P}{S} \right\rfloor + 1
\]
$i$-р цонхны индексүүдийн олонлог: $\mathcal{R}_i = \{\,iS,\; iS+1,\;\dots,\; iS+P-1\,\}$

\subsection*{4.2 Урагшлах тархалт (ногоон хайрцаг)}
\begin{equation}
\boxed{\;\xout[i] \;=\; \max_{k \in \mathcal{R}_i} \xin[k], \qquad i = 0,\dots,M-1\;}
\end{equation}
Урагшлах үед \textbf{ялагчийн индексийг хадгалах} ёстой:
\begin{equation}
k^{*}(i) \;=\; \arg\max_{k \in \mathcal{R}_i} \xin[k]
\end{equation}
Энэ индексийг ``switch'' эсвэл ``argmax mask'' гэж нэрлэдэг ба буцах тархалтад
зайлшгүй хэрэгтэй.

\subsection*{4.3 Буцах тархалт — оролт руу (улаан хайрцаг)}
Зөвхөн ялагч элемент гаралтад нөлөөлсөн тул:
\[
\frac{\partial\, \xout[i]}{\partial\, \xin[j]} =
\begin{cases} 1, & j = k^{*}(i) \\ 0, & \text{бусад} \end{cases}
\]
Иймээс:
\begin{equation}
\boxed{\;\dL{\xin[j]} \;=\; \sum_{\;i\,:\;k^{*}(i)\,=\,j} \dL{\xout[i]}\;}
\end{equation}
\begin{itemize}[leftmargin=*]
  \item \textbf{Давхцалгүй цонх} ($S \ge P$): индекс бүр хамгийн ихдээ \emph{нэг} цонхны
        ялагч байна $\Rightarrow$ градиент зүгээр л \emph{чиглүүлэгдэнэ} (routing),
        нийлбэр хэрэггүй.
  \item \textbf{Давхцалтай цонх} ($S < P$): нэг элемент хэд хэдэн цонхны ялагч байж
        болно $\Rightarrow$ градиентүүдийг \textbf{нэмнэ} (дахин \emph{branch} дүрэм).
\end{itemize}

\subsection*{4.4 Жингийн градиент (гурав дахь хайрцаг)}
\begin{equation}
\boxed{\;\text{Max pooling нь \textbf{параметргүй} тул } \dL{\mathbf{W}} \text{ оршин байхгүй.}\;}
\end{equation}
Лекцийн слайд дээрх жин шинэчлэх хайрцаг энд \textbf{хоосон} байна. Max pooling нь
сурдаг зүйлгүй, зөвхөн \emph{тогтмол} үйлдэл гүйцэтгэдэг.

\subsection*{4.5 Дифференциалагдах чанар}
$\max$ функц нь \textbf{тэнцүү утгууд дээр} дифференциалагдахгүй (жишээ нь
$\xin = (3,3)$). Практикт:
\begin{itemize}[leftmargin=*]
  \item хамгийн \emph{бага индекстэй} ялагчийг сонгодог (NumPy/PyTorch-ийн анхдагч), эсвэл
  \item градиентийг тэнцүү ялагчдад хуваадаг.
\end{itemize}
Тэнцэх магадлал нь хэмжээс тэг (measure zero) олонлог тул бодит өгөгдөлд нөлөөлөхгүй.

\subsection*{4.6 Тоон жишээ}
$\xin = (1,\,-3,\,5,\,2,\,-1,\,4)$, \quad $P = S = 2$, \quad $N=6 \Rightarrow M=3$.

\begin{center}
\begin{tabular}{clcc}
\toprule
$i$ & цонх $\mathcal{R}_i$ & $\xout[i]$ & $k^{*}(i)$ \\
\midrule
0 & $\{1,\,-3\}$  & 1 & 0 \\
1 & $\{5,\,2\}$   & 5 & 2 \\
2 & $\{-1,\,4\}$  & 4 & 5 \\
\bottomrule
\end{tabular}
\end{center}
$\Longrightarrow \xout = (1,\,5,\,4)$

\medskip
$\dL{\xout} = (0.5,\,-1,\,2)$ гэж авбал градиент зөвхөн $k^{*}$ байрлалууд руу очно:
\[
\dL{\xin} = (\,\underbrace{0.5}_{j=0},\; 0,\; \underbrace{-1}_{j=2},\; 0,\; 0,\; \underbrace{2}_{j=5}\,)
\]
Тэг биш элементийн тоо нь яг \textbf{цонхны тоотой тэнцүү} (3) байгааг анзаар.

\subsection*{4.7 Хил хязгаарын тухай хэлэлцүүлэг}
$(N-P)$ нь $S$-д бүхлээр хуваагдахгүй бол сүүлийн цонх \textbf{бүрэн бус} болно.
Гурван арга бий:
\begin{enumerate}[leftmargin=*]
  \item \textbf{Хаях} (floor, PyTorch-ийн \texttt{ceil\_mode=False} анхдагч).
        Сүүлийн хэдэн элемент огт уншигдахгүй $\Rightarrow$ тэдний градиент $0$.
        Хамгийн энгийн бөгөөд аюулгүй.
  \item \textbf{Дүүргэх.} Энд \textbf{маш чухал нарийн ширийн зүйл} байна:
  \begin{mdframed}[backgroundcolor=red!8]
  Max pooling-ийг \textbf{тэгээр дүүргэж БОЛОХГҮЙ}. Хэрэв бүх бодит утга сөрөг байвал
  ($\xin = (-5,-2)$ гэх мэт) зохиомол $0$ нь максимумыг \emph{хожиж}, градиент нь
  \emph{байхгүй} байрлал руу очно. \\[4pt]
  Зөв нь $-\infty$ (эсвэл тухайн төрлийн хамгийн бага утга) -аар дүүргэх:
  \[
  \max(a,\,-\infty) = a \qquad \text{яг л} \qquad a + 0 = a
  \]
  \textbf{Гүн шалтгаан:} $0$ нь \emph{нийлбэрийн} нейтраль элемент, $-\infty$ нь
  \emph{максимумын} нейтраль элемент. Конволюцид нийлбэр хийдэг тул $0$ зөв;
  pooling-д максимум авдаг тул $-\infty$ зөв.
  \end{mdframed}
  \item \textbf{Бүрэн бус цонхыг хэсэгчлэн авах} (\texttt{ceil\_mode=True}).
        Цонхыг оруулна, гэхдээ зөвхөн \emph{хүчинтэй} элементүүд дээр максимум авна.
\end{enumerate}
\end{solution}

%==================================================================================
\section{Бодлого 5 — Backpropagation, нэг алхам}
%==================================================================================

\begin{problem}{5 (слайд 139, бодолт нь 141--147)}
Дараах сүлжээний \textbf{нэг алхмын буцах тархалтыг тооцоол}:
\[
\mathbf{x}_0 \;\to\; \boxed{\mathbf{W}_0\mathbf{x}_0} \;\to\; \mathbf{x}_1 \;\to\;
\boxed{\tanh(\mathbf{x}_1)} \;\to\; \mathbf{x}_2 \;\to\;
\boxed{\mathbf{W}_1\mathbf{x}_2} \;\to\; \mathbf{x}_3 \;\to\;
\boxed{\mathcal{L} = \tfrac12\|\mathbf{x}_3-\mathbf{y}\|_2^2}
\]
\[
\mathbf{W}_0 = \begin{pmatrix} 1 & -3 \\ 0.2 & 1 \end{pmatrix},\quad
\mathbf{W}_1 = \begin{pmatrix} 1 & -1 \end{pmatrix},\quad
\mathbf{x}_0 = \begin{pmatrix} 1.0 \\ 0.1 \end{pmatrix},\quad
\mathbf{y} = 0.5,\quad \eta = 0.2
\]
\end{problem}

\begin{solution}
\subsection*{5.1 Хэрэглэх томъёонууд}
\begin{center}
\begin{tabular}{lccc}
\toprule
\textbf{Модуль} & \textbf{Forward} & \textbf{Backward (оролт руу)} & \textbf{Жингийн градиент} \\
\midrule
Linear    & $\xout = \mathbf{W}\xin$ & $\dL{\xin}=\dL{\xout}\mathbf{W}$ & $\dL{\mathbf{W}}=\xin\dL{\xout}$ \\[3pt]
Pointwise & $\xout = h(\xin)$        & $\dL{\xout}\odot h'(\xin)$       & --- \\[3pt]
$L_2$ loss& $\tfrac12\|\xin-\mathbf{y}\|^2$ & $(\xin-\mathbf{y})^T$     & --- \\
\bottomrule
\end{tabular}
\end{center}
Хэрэгтэй уламжлал: $\;\tanh'(z) = 1 - \tanh^2(z)$

\subsection*{5.2 Forward pass}
\begin{align}
\mathbf{x}_1 &= \mathbf{W}_0\mathbf{x}_0
= \begin{pmatrix} 1 & -3 \\ 0.2 & 1 \end{pmatrix}
  \begin{pmatrix} 1.0 \\ 0.1 \end{pmatrix}
= \begin{pmatrix} 1\cdot1.0 + (-3)\cdot 0.1 \\ 0.2\cdot 1.0 + 1\cdot 0.1 \end{pmatrix}
= \begin{pmatrix} 0.7 \\ 0.3 \end{pmatrix} \\[6pt]
\mathbf{x}_2 &= \tanh(\mathbf{x}_1)
= \begin{pmatrix} \tanh(0.7) \\ \tanh(0.3) \end{pmatrix}
= \begin{pmatrix} 0.604368 \\ 0.291313 \end{pmatrix} \\[6pt]
\mathbf{x}_3 &= \mathbf{W}_1\mathbf{x}_2
= 1\cdot 0.604368 + (-1)\cdot 0.291313 = 0.313055 \\[6pt]
\mathcal{L} &= \tfrac12(\mathbf{x}_3 - \mathbf{y})^2
= \tfrac12(0.313055 - 0.5)^2 = \tfrac12(-0.186945)^2 = \mathbf{0.017474}
\end{align}

\subsection*{5.3 Backward pass}
\paragraph{(i) Алдагдлын давхарга.}
\begin{equation}
\dL{\mathbf{x}_3} = \mathbf{x}_3 - \mathbf{y} = 0.313055 - 0.5 = \mathbf{-0.186945}
\end{equation}
Сөрөг тэмдэг нь ``$\mathbf{x}_3$-ыг \emph{өсгөх} хэрэгтэй'' гэсэн утгатай.

\paragraph{(ii) Хоёр дахь шугаман давхаргаар буцах.}
\begin{equation}
\dL{\mathbf{x}_2} = \dL{\mathbf{x}_3}\,\mathbf{W}_1
= -0.186945\begin{pmatrix} 1 & -1 \end{pmatrix}
= \begin{pmatrix} -0.186945 & 0.186945 \end{pmatrix}
\end{equation}

\paragraph{(iii) $\tanh$ давхаргаар буцах.}
$\tanh$ нь \textbf{цэг бүрийн (pointwise)} үйлдэл тул түүний Jacobian нь
\textbf{диагональ матриц} болно:
\[
\frac{\partial \tanh(\mathbf{x}_1)}{\partial \mathbf{x}_1} =
\begin{pmatrix} 1-\tanh^2(0.7) & 0 \\ 0 & 1-\tanh^2(0.3) \end{pmatrix}
= \begin{pmatrix} 0.634740 & 0 \\ 0 & 0.915137 \end{pmatrix}
\]
Диагональ матрицаар үржүүлэх нь \emph{гишүүнчлэн} үржүүлэхтэй адил:
\begin{equation}
\dL{\mathbf{x}_1}
= \begin{pmatrix} -0.186945\cdot 0.634740 & 0.186945\cdot 0.915137 \end{pmatrix}
= \begin{pmatrix} \mathbf{-0.118661} & \mathbf{0.171080} \end{pmatrix}
\end{equation}

\subsection*{5.4 Жингийн градиентууд}
\begin{align}
\dL{\mathbf{W}_0} &= \mathbf{x}_0 \,\dL{\mathbf{x}_1}
= \begin{pmatrix} 1.0 \\ 0.1 \end{pmatrix}
  \begin{pmatrix} -0.118661 & 0.171080 \end{pmatrix}
= \begin{pmatrix} -0.118661 & 0.171080 \\ -0.011866 & 0.017108 \end{pmatrix} \\[8pt]
\dL{\mathbf{W}_1} &= \mathbf{x}_2 \,\dL{\mathbf{x}_3}
= \begin{pmatrix} 0.604368 \\ 0.291313 \end{pmatrix}(-0.186945)
= \begin{pmatrix} -0.112983 \\ -0.054459 \end{pmatrix}
\end{align}

\begin{mdframed}[backgroundcolor=red!8]
\textbf{Слайд 145 дээрх алдаа.} Лекцийн слайдад
$\dL{\mathbf{W}_1} = \mathbf{x}_2 \cdot (-0.1186)$ гэж бичсэн байна. Энэ нь \textbf{буруу}:
$-0.1186$ нь $\dL{\mathbf{x}_1}$-ийн эхний гишүүн, харин энд $\dL{\mathbf{x}_3} = -0.186945$
хэрэгтэй. \\[4pt]
\textbf{Шалгалт:} $0.604368 \times (-0.186945) = -0.112983$ \checkmark{} (слайдын
$-0.113$ хариутай таарч байна), харин $0.604368 \times (-0.1186) = -0.071678$ \ding{55}.
Өөрөөр хэлбэл \emph{хариу нь зөв, бичсэн тоо нь буруу} — слайдын хэвлэлийн алдаа.
\end{mdframed}

\subsection*{5.5 Жин шинэчлэх}
\[
\mathbf{W}^{k+1} = \mathbf{W}^{k} - \eta\left(\dL{\mathbf{W}}\right)^{T}, \qquad \eta = 0.2
\]
\paragraph{$\mathbf{W}_0$-ийн хувьд.} Эхлээд транспоз хийнэ:
\[
\left(\dL{\mathbf{W}_0}\right)^{T} =
\begin{pmatrix} -0.118661 & -0.011866 \\ 0.171080 & 0.017108 \end{pmatrix}
\]
\begin{align}
\mathbf{W}_0^{k+1}
&= \begin{pmatrix} 1 & -3 \\ 0.2 & 1 \end{pmatrix}
 - 0.2\begin{pmatrix} -0.118661 & -0.011866 \\ 0.171080 & 0.017108 \end{pmatrix} \nonumber\\[4pt]
&= \begin{pmatrix} 1 + 0.023732 & -3 + 0.002373 \\ 0.2 - 0.034216 & 1 - 0.003422 \end{pmatrix}
 = \begin{pmatrix} \mathbf{1.0237} & \mathbf{-2.9976} \\ \mathbf{0.1658} & \mathbf{0.9966} \end{pmatrix}
\end{align}

\paragraph{$\mathbf{W}_1$-ийн хувьд.}
\begin{equation}
\mathbf{W}_1^{k+1}
= \begin{pmatrix} 1 & -1 \end{pmatrix}
 - 0.2\begin{pmatrix} -0.112983 & -0.054459 \end{pmatrix}
= \begin{pmatrix} \mathbf{1.0226} & \mathbf{-0.9891} \end{pmatrix}
\end{equation}

Лекцийн слайд 2 оронгоор дугуйлсан хариу:
$\mathbf{W}_0 = \begin{psmallmatrix} 1.02 & -3.0 \\ 0.17 & 1.0 \end{psmallmatrix}$,
$\mathbf{W}_1 = \begin{pmatrix} 1.02 & -0.989 \end{pmatrix}$ — \textbf{бүрэн таарч байна}.

\begin{mdframed}[backgroundcolor=yellow!12]
\textbf{Транспозыг бүү мартаарай.} $\dL{\mathbf{W}_0}$-ийн мөр нь $\mathbf{x}_0$-ийн
индекс, багана нь $\mathbf{x}_1$-ийн индекс; харин $\mathbf{W}_0$-д эсрэгээр.
Транспозыг мартвал $W_0[1,0]$ нь $0.2024$ гарна (зөв нь $\mathbf{0.1658}$) — шалгалтанд
хамгийн элбэг гардаг алдаа.
\end{mdframed}

\subsection*{5.6 Шалгалт: алдагдал буурсан уу?}
Шинэ жингээр forward pass-ыг дахин гүйцэтгэвэл:
\begin{align*}
\mathbf{x}_1' &= (0.723969,\; 0.265784) \\
\mathbf{x}_2' &= (0.619361,\; 0.259717) \\
\mathbf{x}_3' &= 0.376804 \\
\mathcal{L}'  &= \tfrac12(0.376804-0.5)^2 = \mathbf{0.007589}
\end{align*}
\[
\boxed{\;\mathcal{L}: \; 0.017474 \;\longrightarrow\; 0.007589
\qquad (\text{\textbf{56.6\%} буурсан})\;}
\]
Мөн $\mathbf{x}_3$ нь $0.3131 \to 0.3768$ болж зорилтот $\mathbf{y}=0.5$ руу
\emph{ойртсон}. Градиентийн алхам зөв чиглэлд явсан нь баталгаажлаа. $\blacksquare$
\end{solution}

%==================================================================================
\appendix
\section{Хавсралт А — Тоон баталгаажуулалт}
%==================================================================================

Бүх бодолтыг NumPy дээр бие даан шалгав. Конволюци ба max pooling-ийн градиентийг
\textbf{тоон градиентийн шалгалтаар} (central difference, $\varepsilon = 10^{-6}$)
баталсан.

\begin{center}
\begin{tabular}{llr}
\toprule
\textbf{Бодлого} & \textbf{Шалгасан зүйл} & \textbf{Хамгийн их алдаа} \\
\midrule
3 (Conv)  & $\partial\mathcal{L}/\partial \xin$ & $2.77\times10^{-10}$ \\
3 (Conv)  & $\partial\mathcal{L}/\partial \mathbf{W}$ & $2.10\times10^{-10}$ \\
4 (Pool)  & $\partial\mathcal{L}/\partial \xin$ & $9.84\times10^{-11}$ \\
5 (Backprop) & Бүх завсрын утга & аналитиктай бүрэн тэнцүү \\
\bottomrule
\end{tabular}
\end{center}

\subsection*{A.1 Конволюцийн модулийн шалгалт}
\begin{codebox}\begin{verbatim}
import numpy as np
N, K = 8, 3

def conv_fwd(x, W):            # x_out[i] = sum_k W[k]*x_in[i+k]
    xp = np.concatenate([x, np.zeros(K-1)])    # тэгээр дүүргэх
    return np.array([sum(W[k]*xp[i+k] for k in range(K)) for i in range(N)])

def conv_bwd(x, W, g):         # g = dL/dx_out
    gW = np.array([sum(g[i]*(x[i+k] if i+k < N else 0.0)
                       for i in range(N)) for k in range(K)])
    gx = np.array([sum(g[i]*W[j-i] for i in range(N) if 0 <= j-i < K)
                   for j in range(N)])
    return gx, gW
# Тоон градиенттэй харьцуулахад алдаа ~1e-10
\end{verbatim}\end{codebox}

\subsection*{A.2 Max pooling модулийн шалгалт}
\begin{codebox}\begin{verbatim}
P = S = 2

def pool_fwd(x):
    return np.array([x[i*S:i*S+P].max() for i in range(N//S)])

def pool_bwd(x, g):
    gx = np.zeros(N)
    for i in range(N//S):
        gx[i*S + np.argmax(x[i*S:i*S+P])] += g[i]   # += нь давхцалтай цонхонд чухал
    return gx
\end{verbatim}\end{codebox}

\subsection*{A.3 Backpropagation-ы шалгалт}
\begin{codebox}\begin{verbatim}
W0 = np.array([[1.0, -3.0], [0.2, 1.0]]);  W1 = np.array([[1.0, -1.0]])
x0 = np.array([[1.0], [0.1]]);  y = 0.5;  eta = 0.2

x1 = W0 @ x0;  x2 = np.tanh(x1);  x3 = (W1 @ x2)[0, 0]
L  = 0.5*(x3 - y)**2                          # 0.017474

g3 = x3 - y                                   # -0.186945
g2 = g3 * W1                                  # (-0.186945, 0.186945)
g1 = g2 * (1 - np.tanh(x1.ravel())**2)        # (-0.118661, 0.171080)

gW0 = x0 @ g1;   gW1 = x2 * g3
W0n = W0 - eta*gW0.T;   W1n = W1 - eta*gW1.T  # транспозыг анхаар!
# Шинэ алдагдал: 0.007589  (56.6% буурсан)
\end{verbatim}\end{codebox}

\vfill
\hrule
\smallskip
\noindent\footnotesize
\textbf{Эх сурвалж:} F.CS308 Машин сургалт, Лекц 5 — Neural Networks (147 слайд).
Нэмэлт уншлага: Bengio, Courville, Goodfellow, \emph{Deep Learning},
\url{https://www.deeplearningbook.org/contents/mlp.html}

\end{document}
